% ─── Chapter 10: Conclusions & Next Steps ──────────────────────────
\chapter{Conclusions \& Next Steps}
\label{ch:conclusions}

\section{Summary of Achievements}

Over the course of this project, we have built a complete, reproducible RL
training system for Pok\'{e}mon battles:

\begin{itemize}
  \item \textbf{End-to-end training pipeline:} From Showdown server management
        through distributed PPO training to MLflow-tracked evaluation, the
        system requires only a single command to start a training run.
  \item \textbf{Distributed infrastructure:} 48~parallel environments across
        8~Showdown servers, coordinated by Ray RLlib with custom environment
        connectors.
  \item \textbf{Transformer-based policy:} A 10M-parameter model with
        vocabulary embeddings, positional and role encodings, cross-team
        attention bias, and compressed 14-action space.
  \item \textbf{Systematic debugging:} Three observation pipeline bugs found
        and fixed (weather, switch alignment, mask contamination).  Attention
        collapse in deep transformers diagnosed and resolved through
        architectural changes.
  \item \textbf{Validation suite:} Fixed-protocol evaluation with frozen team
        manifests, per-opponent-type metrics, and MLflow-integrated reporting.
  \item \textbf{Self-play infrastructure:} Hot-reloading opponent that tracks
        the training agent's improving policy.
\end{itemize}

\section{Lessons Learned}

\begin{enumerate}
  \item \textbf{Systematic debugging is essential.}  Three bugs in the
        observation pipeline collectively affected 14\,\% of training steps.
        Without a pipeline validator and step-level logging, these would have
        been nearly impossible to identify.

  \item \textbf{Attention collapse is real.}  In 4-layer post-norm
        transformers, deeper layers become dead weight.  Measuring attention
        distributions---not just loss curves---is necessary to diagnose whether
        the model is actually using its capacity.

  \item \textbf{Reward design outweighs architecture tuning.}  The biggest
        performance gains came from fixing observation bugs (win rate vs.\
        random: 65\,\% $\to$ 85\,\%), not from architectural changes.  Once
        the architecture was no longer the bottleneck, further architecture
        changes produced no win-rate improvement---the reward signal became
        the binding constraint.

  \item \textbf{Curriculum learning provides structure but can mask problems.}
        The easy stage was solved quickly, giving a false sense of progress.
        The real challenge only appeared when facing heuristic opponents.

  \item \textbf{Metrics must be tracked by opponent type.}  An aggregate
        ``50\,\% win rate'' obscured the fact that the agent was winning all
        random games and losing all heuristic games.
\end{enumerate}

\section{What Did Not Work}

\begin{itemize}
  \item \textbf{4-layer post-norm transformer:} Layers 2--3 were dead (uniform
        attention).  Effective depth: 2 out of 4 layers.
  \item \textbf{Hash-based embeddings:} 7.5\,\% collision rate degraded entity
        distinction.
  \item \textbf{22-action native space:} Large branching factor slowed
        exploration.  Compressed to 14 actions without loss of coverage.
  \item \textbf{Disabled LSTM:} Without cross-turn memory, the agent treated
        each turn independently, unable to track opponent move patterns or
        plan multi-turn sequences.
\end{itemize}

\section{Actionable Next Steps}

\begin{enumerate}
  \item \textbf{Reward signal redesign} (highest priority): Experiment with
        auxiliary tasks (opponent-move prediction, state-value prediction as
        a supervised objective), more granular shaped rewards, and reward
        normalisation to improve value-function learning.

  \item \textbf{Complete gauntlet validation pipeline:} Implement automated
        evaluation against the full BDSP gauntlet sequence with fixed seeds,
        deterministic configuration, and standardised metric outputs.

  \item \textbf{Validate LSTM in full training:} Run a controlled experiment
        comparing LSTM-enabled vs.\ LSTM-disabled training to quantify the
        benefit of cross-turn memory.

  \item \textbf{Achieve Milestone~1 competency:} Demonstrate high win rate
        against heuristic opponents under the validation protocol.

  \item \textbf{Begin Milestone~2 (random teams):} Extend to randomly
        generated teams for both agent and opponent, testing generalisation
        under high entropy.

  \item \textbf{Team builder research (Milestone~3):} Investigate approaches
        for a team composition agent---likely requiring a separate model or a
        hierarchical policy structure.
\end{enumerate}

\section{Milestone Outlook}

\begin{center}
\small
\begin{tabularx}{\textwidth}{c l X}
  \toprule
  \textbf{Milestone} & \textbf{Status} & \textbf{Path Forward} \\
  \midrule
  M1: Proof of Concept & In progress &
    Reward redesign + validation pipeline.  Expected resolution within
    2--4 weeks. \\
  M2: Random Battles & Planned &
    Extends M1 with random team generation.  Estimated 4--6 weeks
    after M1. \\
  M3: Team Building & Future &
    Requires team-composition model.  Estimated 10--13 weeks total.
    Largest technical risk. \\
  M4: Gauntlet Master & Future &
    Full Nuzlocke gauntlet.  Depends on M1--M3.  Estimated 12--18 weeks
    total. \\
  \bottomrule
\end{tabularx}
\end{center}

The project has established a solid engineering foundation.  The path from
current state to Milestone~1 completion is clear and primarily requires
iterative reward-signal experimentation rather than further architectural
changes.
